% music-chords.tex — chord symbols (\MusicChord) and Roman numerals.

\providecommand{\RomNum}[1]{\mbox{\ensuremath{\mathrm{#1}\!}}}
\providecommand{\RomNumS}[2]{\mbox{\ensuremath{\mathrm{#1}^{#2}\!}}}
\providecommand{\RomViiO}{\mbox{\ensuremath{\mathrm{vii}^\circ}}}

\newcommand{\musicchordLetterAcc}[1]{\mbox{\raisebox{0.45ex}{\ensuremath{#1}}}}

% \MusicChord keys:
%   root, acc (flat|sharp), quality (maj|min|dom|aug), seven, alts, bass, bassacc
%   slash, slashacc — aliases for bass / bassacc
% Example:
%   \MusicChord{root=C, quality=maj, seven=7, alts={9/flat,5/flat}, bass=E, bassacc=flat}
% Each alts item is degree/kind with kind one of: flat, sharp (e.g. 11/sharp).

\directlua{
  local loader_path = kpse.find_file("source/lua/book_loader.lua")
    or "source/lua/book_loader.lua"

  if loader_path:sub(1, 1) ~= "/" and lfs and lfs.currentdir then
    loader_path = lfs.currentdir() .. "/" .. loader_path
  end

  bookloader = dofile(loader_path)
}

\NewDocumentCommand \MusicChord { m }
  {
    \directlua{
      bookloader.call("music_chords", "print", "\luaescapestring{#1}")
    }
  }
